\section{Discussion}
\subsection{(what determines transferability?)}

Taken together, our results across three tasks suggest that both the direction and the strength of transfer depend on the nature of the visual attributes a model must internalize to solve each task.


For counting and spatial relations, the key attributes are coarse semantic signals such as object numerosity and layout. These signals tend to be encoded at a similar level regardless of whether the model is trained with an understanding-oriented or generation-oriented objective. Notably, the semantic representations learned through understanding can also provide effective guidance for generation. This observation aligns with prior work showing that improving semantic representations can directly benefit generative performance~\cite{yu2024repa, zheng2025rae}. In our experiments, a model trained to understand counting develops representations that reliably capture the number of objects in an image, and these representations remain useful when the model is asked to generate images containing a specified number of objects. Similarly, training for spatial-relation understanding encourages representations that encode relative object configurations, which can be exploited during generation to produce the intended layout. In other words, semantic knowledge acquired through understanding can meaningfully condition the generation process. When the shared concept is sufficiently coarse and can be captured at a semantic level, this conditioning effect can operate in both directions, leading to bi-directional transfer.

In contrast, text recognition and text generation exhibit an asymmetric transfer pattern. A generation objective requires the model to synthesize individual characters with pixel-level fidelity, and the resulting fine-grained character representations transfer naturally to recognition. The reverse direction is comparatively weaker. While prior studies report that recognition training can facilitate text generation~\cite{ji2023improving, fang2025recognsyn, min2025tair}, we also observe only limited transfer from understanding to generation in this setting. A plausible explanation is that recognition-oriented training primarily emphasizes high-level textual content and identity, whereas text generation additionally demands low-level, spatially precise rendering details. Consequently, semantic knowledge extracted from recognition alone is insufficient to fully bridge the gap to pixel-accurate character synthesis. Overall, our findings indicate that when the shared concept requires fine-grained, pixel-level precision, semantic representations from understanding provide only partial support for generation.

\subsection{}
% transferability 의 application potential

